\documentclass[12pt]{article}                                                                                                           
\usepackage[utf8]{inputenc}                                                                                                             
\usepackage[T1]{fontenc}                                                                                                                
\usepackage{amsmath, amssymb}                                                                                                           
\usepackage{geometry}                                                                                                                   
\usepackage{xcolor}                                                                                                                     
\usepackage{lipsum}                                                                                                                     
\usepackage{titlesec}                                                                                                                   
\usepackage{fancyhdr}                                                                                                                   
\usepackage[most]{tcolorbox}                                                                                                                  
\usepackage{graphicx}                                                                                                                   
\usepackage{float}                                                                                      
\usepackage{tikz}
\usetikzlibrary{shapes.geometric, arrows.meta}
\usetikzlibrary{arrows.meta, decorations.markings}                                                                                      
                                                                                                                                        
\geometry{a4paper, margin=2.5cm}                                                                                                        
\pagestyle{fancy}                                                                                                                       
\fancyhf{}                                                                                                                              
\rhead{Problem Set V}                                                                                                                   
\lhead{Rigid Body Mechanics}                                                                                                                     
\cfoot{\thepage}                                                                                                                        
                                                                                                                                        
% Fancy titles                                                                                                                          
\titleformat{\section}{\normalfont\Large\bfseries}{Problem \thesection:}{1em}{}                                                        
\titleformat{\subsection}{\normalfont\bfseries}{Correction:}{1em}{}                                                                     
                                                                                                                                        
% Custom box for correction with breakable option
\newtcolorbox{correctionbox}{                                                                                                           
  colback=gray!5,                                                                                                                       
  colframe=black,                                                                                                                       
  fonttitle=\bfseries,                                                                                                                  
  title=Correction,                                                                                                                     
  before skip=10pt,                                                                                                                     
  after skip=10pt,
  breakable,
  enhanced
}                                                                                                                                       
                                                                                                                                        
% Fix headheight warning
\setlength{\headheight}{14.49998pt}

\begin{document}

\begin{center}
\Large\textbf{Ibn Tofail University} \\[1em]
\large\textit{Rigid Body Mechanics} \\[0.5em]
\large\textit{Problem Set V} \\[0.5em]
\end{center}

\vspace{1cm}

% ----------- COURSE QUESTIONS -------------
\section*{Course Questions*}

1) Show the angular momentum theorem at a mobile point A of the solid in $R_o$
$$\left[\frac{d}{dt}\vec{L}_A(S/R_o)\right]_{R_o} = \vec{M}_A(\Sigma \vec{F}_{ext}) + m\vec{V}(G/R_o) \times \vec{V}(A/R_o)$$

2) Show that the power of the resultant forces acting on a solid is equal to the product (comoment) of the kinematic torsor and the force torsor.
$$\mathcal{P}(S/R_o) = [T_V(S/R_o)]_A \otimes [T_F(S/R_o)]_A = \left\{\begin{array}{c} \vec{\Omega}(S/R_o) \\ \vec{V}(A \in S/R_o) \end{array}\right\} \otimes \left\{\begin{array}{c} \vec{F} \\ \vec{M}_A(\vec{F}) \end{array}\right\}$$

3) a) Write the statement of the kinetic energy theorem.

\quad b) Prove this theorem.

\begin{correctionbox}

\end{correctionbox}

% ----------- PROBLEM 1 -------------
\section{}

A heavy pendulum consists of a homogeneous solid of arbitrary shape, of mass m, rotating around a fixed point O (Figure 1).

This pendulum is fixed at point O by a perfect pivot joint. The pendulum is linked to the reference frame $R_1(O, \vec{x}_1, \vec{y}_1, \vec{z}_1)$ in rotational motion with respect to the fixed reference frame $R_o(O, \vec{x}_o, \vec{y}_o, \vec{z}_o)$. The motion takes place in the plane $(\vec{x}_o, O, \vec{y}_o)$.

Given: $\overrightarrow{OG} = L\vec{x}_1$, $I_G$ the moment of the solid with respect to the axis $(G, \vec{z}_1)$ and $\vec{R} = R_x\vec{x}_1 + R_y\vec{y}_1$ the reaction at point O.

1) Calculate the absolute velocity and acceleration of the center of mass G.

2) By applying the theorems of dynamic resultant and dynamic moment, establish the differential equation of motion.

3) Find the differential equation of motion by applying the kinetic energy theorem.

4) Show that the total mechanical energy is conserved and find the differential equation of motion of the heavy pendulum.

5) Deduce the differential equation of the simple pendulum as well as its period for small oscillations.

\begin{correctionbox}

\end{correctionbox}

% ----------- PROBLEM 2 -------------
\section{}

A solid of revolution S, solid and homogeneous, consists of a hemisphere and a cylinder with the same base. We denote by a the radius of this base and H its center. We note m the mass of S, G its center of inertia which we locate in S by $\overrightarrow{HG} = h\vec{z}$ (figure 2), A and C its principal moments of inertia at G [C moment of inertia relative to $(G, \vec{z})$].

S is in contact, by its hemispherical part only, at a point I with a horizontal plane $(\Pi)$ to which is linked the reference frame $R_o(O, \vec{x}_o, \vec{y}_o, \vec{z}_o)$, $\vec{z}_o$ vertical upward and on the side of S [i.e. $\overrightarrow{OG} \cdot \vec{z}_o > 0$]. The direct orthonormal reference frame $R(G, \vec{x}, \vec{y}, \vec{z})$ is linked to S. We locate the position of S in $R_o$ by the usual Euler angles $\psi$, $\theta$, $\phi$ and by the coordinates x and y in $R_o$ of the orthogonal projection of G on the plane $(\Pi)$.

We note $\vec{z}_o \times \vec{u} = \vec{v}$ and $\vec{z} \times \vec{u} = \vec{w}$ with $\psi = (\vec{x}_o, \vec{u})$.

\textbf{A) Kinematics and kinetics.}

1) Determine the instantaneous rotation velocity of the solid $\vec{\Omega}(S/R_o)$.

2) Calculate the z-coordinate of G in $R_o$.

3) Calculate the velocities $\vec{V}(G/R_o)$ and $\vec{V}(I \in S/R_o)$. Express these velocities in the basis $(\vec{u}, \vec{v}, \vec{z}_o)$. Does the last velocity have a particular meaning?

4) Calculate in the basis $(\vec{u}, \vec{w}, \vec{z})$, the angular momentum of the solid (S) at point G.

5) Calculate the kinetic energy of the solid with respect to the reference frame $R_o(O, \vec{x}_o, \vec{y}_o, \vec{z}_o)$.

\textbf{B) Dynamics.}

The external actions exerted on S are:
\begin{itemize}
\item The uniform and constant gravitational field of acceleration $\vec{g} = -g\vec{z}_o$.
\item The reaction of the plane $(\Pi)$ on the solid S defined by a slider whose support passes through I and of vector $\vec{R} = R\vec{z}_o$ (no friction).
\item The action of an unspecified device, defined by a slider whose support passes through G and of vector $\vec{F} = -k\overrightarrow{OG}$ (k positive constant given).
\end{itemize}

1) By application of the dynamic resultant theorem calculate explicitly x(t) and y(t).

2) By applying the angular momentum theorem at point G, show that:
$$\vec{L}_G(S/R_o) \cdot \vec{z}_o = \text{cte} = k_1 \quad \text{and} \quad \vec{L}_G(S/R_o) \cdot \vec{z} = \text{cte} = k_2$$
Deduce two equations of motion in the form of first integrals.

3) By application of the kinetic energy theorem, give a third equation of motion in the form of a first integral.

\begin{correctionbox}

\end{correctionbox}

% ----------- PROBLEM 3 -------------
\section{}

A homogeneous bar of length $AB = 2L$, of center G and mass m, slides without friction along a staircase as shown in the figure. Point A slides on the ground and point C slides on the edge of the staircase. The initial position of the bar being $A_oB_o$.

$R_1(A, \vec{x}_1, \vec{y}_1, \vec{z}_1)$ reference frame linked to the bar such that $\vec{z}_1 = \vec{z}_0$.

Given $OA = x(t)$ and $\alpha = (\vec{x}_o, \vec{x}_1) = (\vec{y}_o, \vec{y}_1)$.

The inertia matrix of the bar at G in $R_1$ is given by:
$$I_G(S) = \begin{pmatrix} \frac{mL^2}{3} & 0 & 0 \\ 0 & 0 & 0 \\ 0 & 0 & \frac{mL^2}{3} \end{pmatrix}_{R_1}$$

We will take the reference frame $\mathbf{R_o}(O, \vec{x}_o, \vec{y}_o, \vec{z}_o)$ as the projection reference frame.

1) Determine the position vector $\overrightarrow{OG}$.

2) Calculate the velocity vector $\vec{V}(G/R_o)$ and acceleration $\vec{\gamma}(G/R_o)$.

3) By applying the dynamic resultant theorem to the bar, show the following equations:
$$R_C \cos\alpha = m\ddot{x} - mL(\ddot{\alpha}\cos\alpha - \dot{\alpha}^2\sin\alpha) \quad (1)$$
$$R_A + R_C\sin\alpha = mg - mL(\ddot{\alpha}\sin\alpha + \dot{\alpha}^2\cos\alpha) \quad (2)$$

4) By applying the dynamic moment theorem to the bar at point G, show the following equation:
$$LR_A\sin\alpha - R_C\left(\frac{x}{\sin\alpha} - L\right) = m\frac{L^2}{3}\ddot{\alpha} \quad (3)$$

5) By applying the kinetic energy theorem to the bar, show that the mechanical energy of the bar is conserved $(E_m = \text{cte})$. Deduce the first integral of energy.
$$\frac{L^2}{3}\dot{\alpha}^2 + \dot{x}^2 + L^2\dot{\alpha}^2 - 2L\dot{x}\dot{\alpha}\cos\alpha + 2gL\cos\alpha = \text{cte} \quad (4)$$

\begin{correctionbox}

\end{correctionbox}

\end{document}